%%%%%%%%%%%%%%%%%%%%%%%%%%%%%%%%%%%%%%%%%%%%%%%%%%%%%%%%%%%%%%%%%%%%%%%%%%%%%%%
% Harvard Academic Notes - 통합 마스터 템플릿
% 모든 강의 노트에 적용되는 통일된 스타일
% 버전: 2.1 - 가독성 개선 (선택적 최적화)
% 최종 수정일: 2025-11-17
%%%%%%%%%%%%%%%%%%%%%%%%%%%%%%%%%%%%%%%%%%%%%%%%%%%%%%%%%%%%%%%%%%%%%%%%%%%%%%%

\documentclass[11pt,a4paper]{article}

%========================================================================================
% 기본 패키지
%========================================================================================

% --- 한국어 지원 ---
\usepackage{kotex}

% --- 페이지 레이아웃 ---
\usepackage[top=20mm, bottom=20mm, left=20mm, right=18mm]{geometry}
\usepackage{setspace}
\onehalfspacing                      % 1.5배 줄간격
\setlength{\parskip}{0.5em}          % 문단 간격
\setlength{\parindent}{0pt}          % 들여쓰기 없음

% --- 표 관련 ---
\usepackage{booktabs}              % 고품질 표
\usepackage{tabularx}              % 자동 너비 조절 표
\usepackage{array}                 % 표 컬럼 확장
\usepackage{longtable}             % 여러 페이지 표
\renewcommand{\arraystretch}{1.1}  % 표 행간 조절

%========================================================================================
% 헤더 및 푸터
%========================================================================================

\usepackage{fancyhdr}
\pagestyle{fancy}
\fancyhf{}
\fancyhead[L]{\small\textit{CS109A: Introduction to Data Science}}
\fancyhead[R]{\small\textit{Module 23}}
\fancyfoot[C]{\thepage}
\renewcommand{\headrulewidth}{0.5pt}
\renewcommand{\footrulewidth}{0.3pt}

% 첫 페이지는 헤더 없음
\fancypagestyle{firstpage}{
    \fancyhf{}
    \fancyfoot[C]{\thepage}
    \renewcommand{\headrulewidth}{0pt}
}

%========================================================================================
% 색상 정의 (파스텔 톤 + 다크모드 호환)
%========================================================================================

\usepackage[dvipsnames]{xcolor}

% 밝은 배경용 파스텔 색상
\definecolor{lightblue}{RGB}{220, 235, 255}      % 부드러운 파랑
\definecolor{lightgreen}{RGB}{220, 255, 235}     % 부드러운 초록
\definecolor{lightyellow}{RGB}{255, 250, 220}    % 부드러운 노랑
\definecolor{lightpurple}{RGB}{240, 230, 255}    % 부드러운 보라
\definecolor{lightgray}{gray}{0.95}              % 밝은 회색
\definecolor{lightpink}{RGB}{255, 235, 245}      % 부드러운 핑크
\definecolor{boxgray}{gray}{0.95}
\definecolor{boxblue}{rgb}{0.9, 0.95, 1.0}
\definecolor{boxred}{rgb}{1.0, 0.95, 0.95}
\definecolor{myblue}{RGB}{50, 100, 180}

% 진한 색상 (테두리/제목용)
\definecolor{darkblue}{RGB}{50, 80, 150}
\definecolor{darkgreen}{RGB}{40, 120, 70}
\definecolor{darkorange}{RGB}{200, 100, 30}
\definecolor{darkpurple}{RGB}{100, 60, 150}

%========================================================================================
% 박스 환경 (tcolorbox) - 6가지 타입
%========================================================================================

\usepackage[most]{tcolorbox}
\tcbuselibrary{skins, breakable}

% 1. 개요 박스 (강의 시작 부분)
\newtcolorbox{overviewbox}[1][]{
    enhanced,
    colback=lightpurple,
    colframe=darkpurple,
    fonttitle=\bfseries\large,
    title=📚 강의 개요,
    arc=3mm,
    boxrule=1pt,
    left=8pt,
    right=8pt,
    top=8pt,
    bottom=8pt,
    breakable,
    #1
}

% 2. 요약 박스
\newtcolorbox{summarybox}[1][]{
    enhanced,
    colback=lightblue,
    colframe=darkblue,
    fonttitle=\bfseries,
    title=📝 핵심 요약,
    arc=2mm,
    boxrule=0.7pt,
    left=6pt,
    right=6pt,
    top=6pt,
    bottom=6pt,
    breakable,
    #1
}

% 3. 핵심 정보 박스
\newtcolorbox{infobox}[1][]{
    enhanced,
    colback=lightgreen,
    colframe=darkgreen,
    fonttitle=\bfseries,
    title=💡 핵심 정보,
    arc=2mm,
    boxrule=0.7pt,
    left=6pt,
    right=6pt,
    top=6pt,
    bottom=6pt,
    breakable,
    #1
}

% 4. 주의사항 박스
\newtcolorbox{warningbox}[1][]{
    enhanced,
    colback=lightyellow,
    colframe=darkorange,
    fonttitle=\bfseries,
    title=⚠️ 주의사항,
    arc=2mm,
    boxrule=0.7pt,
    left=6pt,
    right=6pt,
    top=6pt,
    bottom=6pt,
    breakable,
    #1
}

% 5. 예제 박스
\newtcolorbox{examplebox}[1][]{
    enhanced,
    colback=lightgray,
    colframe=black!60,
    fonttitle=\bfseries,
    title=📖 예제,
    arc=2mm,
    boxrule=0.7pt,
    left=6pt,
    right=6pt,
    top=6pt,
    bottom=6pt,
    breakable,
    #1
}

% 6. 정의 박스
\newtcolorbox{definitionbox}[1][]{
    enhanced,
    colback=lightpink,
    colframe=purple!70!black,
    fonttitle=\bfseries,
    title=📌 정의,
    arc=2mm,
    boxrule=0.7pt,
    left=6pt,
    right=6pt,
    top=6pt,
    bottom=6pt,
    breakable,
    #1
}

% 7. 중요 박스 (importantbox - warningbox와 유사)
\newtcolorbox{importantbox}[1][]{
    enhanced,
    colback=boxred,
    colframe=red!70!black,
    fonttitle=\bfseries,
    title=⚠️ 매우 중요,
    arc=2mm,
    boxrule=0.7pt,
    left=6pt,
    right=6pt,
    top=6pt,
    bottom=6pt,
    breakable,
    #1
}

% 8. cautionbox (warningbox와 동일)
\let\cautionbox\warningbox
\let\endcautionbox\endwarningbox

% tcbset 스타일 정의 (\begin{tcolorbox}[summarybox] 형식 사용 시 호환)
\tcbset{
    summarybox/.style={enhanced, colback=lightblue, colframe=darkblue, fonttitle=\bfseries, title=핵심 요약, arc=2mm, boxrule=0.7pt, breakable},
    warningbox/.style={enhanced, colback=lightyellow, colframe=darkorange, fonttitle=\bfseries, title=주의, arc=2mm, boxrule=0.7pt, breakable},
    examplebox/.style={enhanced, colback=lightgray, colframe=black!60, fonttitle=\bfseries, title=예제, arc=2mm, boxrule=0.7pt, breakable},
    definitionbox/.style={enhanced, colback=lightpink, colframe=purple!70!black, fonttitle=\bfseries, title=정의, arc=2mm, boxrule=0.7pt, breakable},
    importantbox/.style={enhanced, colback=boxred, colframe=red!70!black, fonttitle=\bfseries, title=매우 중요, arc=2mm, boxrule=0.7pt, breakable},
    infobox/.style={enhanced, colback=lightgreen, colframe=darkgreen, fonttitle=\bfseries, title=핵심 정보, arc=2mm, boxrule=0.7pt, breakable},
    conceptbox/.style={enhanced, colback=lightgreen, colframe=darkgreen, fonttitle=\bfseries, title=개념, arc=2mm, boxrule=0.7pt, breakable},
    analogybox/.style={enhanced, colback=green!5!white, colframe=green!60!black, fonttitle=\bfseries, title=비유, arc=2mm, boxrule=0.7pt, breakable},
    overviewbox/.style={enhanced, colback=lightpurple, colframe=darkpurple, fonttitle=\bfseries\large, title=강의 개요, arc=3mm, boxrule=1pt, breakable},
    cautionbox/.style={enhanced, colback=lightyellow, colframe=darkorange, fonttitle=\bfseries, title=주의, arc=2mm, boxrule=0.7pt, breakable},
    mybox/.style={enhanced, colback=lightblue, colframe=darkblue, fonttitle=\bfseries, title={#1}, arc=2mm, boxrule=0.7pt, breakable},
    definition/.style={enhanced, colback=lightpink, colframe=purple!70!black, fonttitle=\bfseries, title={#1}, arc=2mm, boxrule=0.7pt, breakable},
    example/.style={enhanced, colback=lightgray, colframe=black!60, fonttitle=\bfseries, title={#1}, arc=2mm, boxrule=0.7pt, breakable},
    summary/.style={enhanced, colback=lightblue, colframe=darkblue, fonttitle=\bfseries, title={#1}, arc=2mm, boxrule=0.7pt, breakable},
    warning/.style={enhanced, colback=lightyellow, colframe=darkorange, fonttitle=\bfseries, title={#1}, arc=2mm, boxrule=0.7pt, breakable},
}

%========================================================================================
% 코드 블록 설정 (밝은 배경)
%========================================================================================

\usepackage{listings}

\definecolor{codegray}{rgb}{0.5,0.5,0.5}
\definecolor{codepurple}{rgb}{0.58,0,0.82}
\definecolor{backcolour}{rgb}{0.95,0.95,0.95}

\lstset{
    basicstyle=\ttfamily\small,
    backgroundcolor=\color{lightgray},
    keywordstyle=\color{darkblue}\bfseries,
    commentstyle=\color{darkgreen}\itshape,
    stringstyle=\color{purple!80!black},
    numberstyle=\tiny\color{black!60},
    numbers=left,
    numbersep=8pt,
    breaklines=true,
    breakatwhitespace=false,
    frame=single,
    frameround=tttt,
    rulecolor=\color{black!30},
    captionpos=b,
    showstringspaces=false,
    tabsize=2,
    xleftmargin=15pt,
    xrightmargin=5pt,
    escapeinside={\%*}{*)}
}

% Python 코드 스타일
\lstdefinestyle{pythonstyle}{
    language=Python,
    morekeywords={self, True, False, None},
}

% SQL 코드 스타일
\lstdefinestyle{sqlstyle}{
    language=SQL,
    morekeywords={SELECT, FROM, WHERE, JOIN, GROUP, BY, ORDER, HAVING},
}

%========================================================================================
% 목차 스타일링
%========================================================================================

\usepackage{tocloft}
\renewcommand{\cftsecleader}{\cftdotfill{\cftdotsep}}
\setlength{\cftbeforesecskip}{0.4em}
\renewcommand{\cftsecfont}{\bfseries}
\renewcommand{\cftsubsecfont}{\normalfont}

%========================================================================================
% 표 및 그림
%========================================================================================

\usepackage{graphicx}              % 이미지
\usepackage{adjustbox}             % 표/박스 크기 조절

% 표 캡션 스타일
\usepackage{caption}
\captionsetup[table]{
    labelfont=bf,
    textfont=it,
    skip=5pt
}
\captionsetup[figure]{
    labelfont=bf,
    textfont=it,
    skip=5pt
}

%========================================================================================
% 수학
%========================================================================================

\usepackage{amsmath, amssymb, amsthm}

% 정리 환경
\theoremstyle{definition}
\newtheorem{theorem}{정리}[section]
\newtheorem{lemma}[theorem]{보조정리}
\newtheorem{proposition}[theorem]{명제}
\newtheorem{corollary}[theorem]{따름정리}
\newtheorem{definition}{정의}[section]
\newtheorem{example}{예제}[section]

%========================================================================================
% 하이퍼링크
%========================================================================================

\usepackage[
    colorlinks=true,
    linkcolor=blue!80!black,
    urlcolor=blue!80!black,
    citecolor=green!60!black,
    bookmarks=true,
    bookmarksnumbered=true,
    pdfborder={0 0 0}
]{hyperref}

% PDF 메타데이터는 각 문서에서 설정
\hypersetup{
    pdftitle={CS109A: Introduction to Data Science - Module 23},
    pdfauthor={강의 노트},
    pdfsubject={Academic Notes}
}

%========================================================================================
% 기타 유용한 패키지
%========================================================================================

\usepackage{enumitem}              % 리스트 커스터마이징
\setlist{nosep, leftmargin=*, itemsep=0.3em}

\usepackage{microtype}             % 타이포그래피 개선
\usepackage{footnote}              % 각주 개선
\usepackage{url}                   % URL 줄바꿈
\urlstyle{same}

%========================================================================================
% 사용자 정의 명령어
%========================================================================================

% 강조 텍스트
\newcommand{\important}[1]{\textbf{\textcolor{red!70!black}{#1}}}
\newcommand{\keyword}[1]{\textbf{#1}}
\newcommand{\term}[1]{\textit{#1}}
\newcommand{\code}[1]{\texttt{#1}}

% 용어 설명 (인라인)
\newcommand{\defterm}[2]{\textbf{#1}\footnote{#2}}

% 섹션 시작 전 페이지 분리
\newcommand{\newsection}[1]{\newpage\section{#1}}

%========================================================================================
% 문서 제목 스타일
%========================================================================================

\usepackage{titling}
\pretitle{\begin{center}\LARGE\bfseries}
\posttitle{\par\end{center}\vskip 0.5em}
\preauthor{\begin{center}\large}
\postauthor{\end{center}}
\predate{\begin{center}\large}
\postdate{\par\end{center}}

%========================================================================================
% 섹션 제목 간격
%========================================================================================

\usepackage{titlesec}
\titlespacing*{\section}{0pt}{1.5em}{0.8em}
\titlespacing*{\subsection}{0pt}{1.2em}{0.6em}
\titlespacing*{\subsubsection}{0pt}{1em}{0.5em}

%========================================================================================
% 메타 정보 박스 명령어
%========================================================================================

\newcommand{\metainfo}[4]{
\begin{tcolorbox}[
    colback=lightpurple,
    colframe=darkpurple,
    boxrule=1pt,
    arc=2mm,
    left=10pt,
    right=10pt,
    top=8pt,
    bottom=8pt
]
\begin{tabular}{@{}rl@{}}
▣ \textbf{강의명:} & #1 \\[0.3em]
▣ \textbf{주차:} & #2 \\[0.3em]
▣ \textbf{교수명:} & #3 \\[0.3em]
▣ \textbf{목적:} & \begin{minipage}[t]{0.75\textwidth}#4\end{minipage}
\end{tabular}
\end{tcolorbox}
}

%========================================================================================
% 끝
%========================================================================================


\begin{document}

\thispagestyle{firstpage}

\metainfo{CS109A}{Lecture 23}{UIUC Student}{본 문서는 CS109A Lecture 23 강의 내용을 바탕으로 부스팅(Boosting) 기법을 심층적으로 다룹니다. 랜덤 포레스트(Random Forest)와의 비교를 통해 부스팅의 동기를 이해하고, 약한 학습기(Weak Learner)를 순차적으로 결합하여 모델의 편향(Bias)을 줄이는 원리를 설명합니다. 특히 그래디언트 부스팅(Gradient Boosting)의 작동 방식을 회귀 예시로 상세히 다루며, 이를 그래디언트 하강(Gradient Descent)의 관점에서 수학적으로 정당화합니다. 표 형식 데이터(Tabular Data)에서 부스팅의 강력한 성능과 학습률(\texttt{\lambda}) 튜닝의 중요성을 강조합니다.}

\tableofcontents

\newpage
\section{용어 정리} % section* -> section으로 변경하여 목차에 포함
다음은 강의에서 다루는 주요 개념들을 쉽게 설명하고 원어를 병기한 표입니다.

\begin{adjustbox}{width=\textwidth,center}
\begin{tabular}{lllc}
\toprule
\textbf{용어} & \textbf{쉬운 설명} & \textbf{원어} & \textbf{비고} \\
\midrule
부스팅 & 여러 약한 모델을 순차적으로 학습시켜 이전 모델의 실수를 보완하며 강력한 모델을 만드는 앙상블 기법. & Boosting & \\
랜덤 포레스트 & 여러 결정 트리를 독립적으로 학습시킨 후 결과를 평균하여 분산을 줄이는 앙상블 기법. & Random Forest & \\
약한 학습기 & 단독으로는 성능이 좋지 않지만, 여러 개를 모으면 강력해지는 간단한 모델. & Weak Learner & 주로 깊이가 얕은 결정 트리(Stump) \\
스텀프 & 깊이가 1인 결정 트리. 즉, 단 한 번의 분할만 수행하는 가장 단순한 트리. & Stump & \\
잔차 & 실제 값과 모델의 예측 값 사이의 차이. 모델이 틀린 정도를 나타냄. & Residual & \\
학습률 & 그래디언트 하강이나 부스팅에서 모델 업데이트 시 얼마나 크게 스텝을 밟을지 결정하는 상수. & Learning Rate (\texttt{\lambda} 또는 \texttt{\eta}) & \\
그래디언트 하강 & 함수의 최솟값을 찾기 위해 기울기(Gradient)의 반대 방향으로 조금씩 이동하는 최적화 알고리즘. & Gradient Descent & \\
편향 & 모델이 실제 관계를 너무 단순하게 가정하여 발생하는 오차. 과소적합(Underfitting)과 관련. & Bias & \\
분산 & 모델이 훈련 데이터의 작은 변화에도 예측이 크게 달라지는 정도. 과대적합(Overfitting)과 관련. & Variance & \\
편향-분산 트레이드오프 & 모델 복잡도를 높이면 편향은 줄지만 분산은 늘고, 그 반대도 마찬가지인 관계. & Bias-Variance Trade-off & \\
표 형식 데이터 & 행과 열로 구성된 테이블 형태의 데이터. 이미지, 텍스트, 오디오와 대비됨. & Tabular Data & \\
과대적합 & 모델이 훈련 데이터에 너무 맞춰져서 새로운 데이터에 대한 예측 성능이 떨어지는 현상. & Overfitting & \\
비모수 모델 & 데이터의 특정 분포를 가정하지 않거나, 고정된 수의 파라미터로 정의되지 않는 모델. & Non-parametric Model & 결정 트리가 대표적 \\
추론 (ML) & 머신러닝 모델이 새로운 데이터에 대해 예측을 수행하는 과정. & Inference (ML) & 통계적 추론과 구분 \\
\bottomrule
\end{tabular}
\caption{강의 주요 용어 정리}
\label{tab:glossary}
\end{adjustbox}

\newpage
\section{부스팅(Boosting) 개요}

\subsection{강의 소개 및 중요성}
오늘 강의는 매우 흥미로운 내용으로 구성되어 있습니다. 처음에는 재미있고 직관적인 개념부터 시작하여, 중간에는 좀 더 구체적인 방법론을 다루고, 마지막에는 수학적인 깊이 있는 내용을 탐구합니다. 온라인으로 시청하는 분들은 현장의 열기를 놓치고 있지만, 이 강의는 모든 것을 담고 있습니다.

특히, 부스팅은 데이터 과학 분야에서 매우 중요한 기법입니다. 이미지, 텍스트, 오디오와 같은 비정형 데이터가 아닌, 행과 열로 잘 정리된 \textbf{표 형식 데이터(Tabular Data)}를 다룰 때, 부스팅과 랜덤 포레스트(Random Forest)는 신경망(Neural Network)보다 더 좋은 성능을 보이거나 최소한 동등한 성능을 냅니다. Kaggle과 같은 데이터 과학 경진대회에서도 표 형식 데이터 부문에서는 부스팅이나 랜덤 포레스트가 거의 항상 우승을 차지합니다. 따라서 이 방법론들을 배우는 것은 매우 중요합니다.

\begin{tipbox}[신경망 vs. 부스팅/랜덤 포레스트] % 마스터 템플릿의 tipbox 사용
신경망이 모든 문제에 대한 만능 해결책은 아닙니다. 특히 표 형식 데이터에서는 부스팅이나 랜덤 포레스트가 더 효율적이고 강력한 성능을 발휘하는 경우가 많습니다. 복잡한 딥러닝 모델이 항상 필요한 것은 아니며, 문제의 특성에 맞는 적절한 도구를 선택하는 것이 중요합니다.
\end{tipbox}

\subsection{이전 학습 내용 복습: 결정 트리와 앙상블}
지금까지 우리는 여러 머신러닝 모델을 학습했습니다.

\begin{itemize}
    \item \textbf{결정 트리(Decision Trees):} 데이터를 분할하여 예측하는 모델입니다.
    \item \textbf{배깅(Bagging):} 부트스트랩(Bootstrap) 샘플링을 통해 여러 트리를 만들고, 이들의 예측을 평균하여 분산(Variance)을 줄이는 방법입니다.
    \item \textbf{랜덤 포레스트(Random Forest):} 배깅의 한 종류로, 각 분할(Split)마다 전체 특성(Predictors) 중 일부를 무작위로 선택하여 트리를 만듭니다. 이는 트리 간의 상관관계를 줄여 다양성(Diversity)을 높이고, 결과적으로 분산을 더욱 효과적으로 감소시킵니다.
\end{itemize}

랜덤 포레스트는 분산 감소에 탁월하지만, 여전히 일반화 성능(Generalization)이 높을 수 있습니다.

\subsection{부스팅의 동기: 편향 감소에 집중}
단일 결정 트리의 단점을 해결하는 다른 방법은 없을까요? 랜덤 포레스트는 깊은 트리를 여러 개 만들어 분산을 줄이는 데 초점을 맞췄습니다. 부스팅은 이와 반대 방향에서 접근합니다.

\begin{tcolorbox}[summary, title={부스팅의 핵심 동기}] % 마스터 템플릿의 summary 스타일 사용
\begin{itemize}
    \item \textbf{얕은 트리(Shallow Trees)의 활용:} 깊이가 얕은 트리는 \textbf{높은 편향(High Bias)}과 \textbf{낮은 분산(Low Variance)}을 가집니다. 즉, 모델이 너무 단순하여 데이터의 복잡한 패턴을 잘 학습하지 못하고(과소적합, Underfitting), 훈련 데이터가 바뀌어도 예측이 크게 변하지 않습니다.
    \item \textbf{편향 감소 목표:} 부스팅은 이러한 얕은 트리의 높은 편향을 줄이는 데 집중합니다.
    \item \textbf{실수로부터 학습:} 이전 모델이 틀린 부분(실수)을 다음 모델이 학습하여 보완하는 방식으로 모델을 점진적으로 개선합니다.
\end{itemize}
\end{tcolorbox}

\begin{figure}[h!]
    \centering
    % \includegraphics[width=0.8\textwidth]{example-bias-variance-tradeoff.png}  % Image not found: example-bias-variance-tradeoff.png % Placeholder for actual image
    \caption{모델 복잡도에 따른 편향-분산 트레이드오프 (가상의 이미지)}
    \label{fig:bias_variance_tradeoff}
\end{figure}

위 그림 \ref{fig:bias_variance_tradeoff}에서 볼 수 있듯이, 모델 복잡도가 증가함에 따라 편향은 감소하고 분산은 증가합니다. 얕은 트리는 복잡도가 낮아 높은 편향과 낮은 분산을 가집니다. 부스팅은 이러한 얕은 트리를 활용하여 편향을 줄이는 방향으로 나아갑니다.

\subsection{직관적 이해: 시험 통과 비유}
부스팅의 아이디어를 쉽게 이해하기 위해 시험 통과 비유를 들어보겠습니다. 여러분은 최종 시험에서 A학점(90점)을 받아야 합니다.

\begin{tcolorbox}[example, title={시험 통과 비유}] % 마스터 템플릿의 example 스타일 사용
\begin{enumerate}
    \item \textbf{첫 번째 약한 학습기:} 친구에게 "시험 잘 보는 팁"을 물어봅니다. 친구는 "절대 D 옵션을 선택하지 마!"라고 조언합니다. 이 규칙은 단순하고 완벽하지 않지만, 시험 점수를 60\%까지 올릴 수 있습니다. (높은 편향, 낮은 분산의 약한 학습기)
    \item \textbf{두 번째 약한 학습기 (실수로부터 학습):} 60점으로는 A학점을 받을 수 없으므로, 이번에는 조교에게 조언을 구합니다. 조교는 "과대적합(Overfitting)이 보기에 있으면 그게 정답일 가능성이 높아!"라고 조언합니다. 이때 중요한 것은, 조교의 조언은 \textbf{이전 규칙으로 틀린 문제에 집중}하여 보완하는 것입니다.
    \item \textbf{세 번째 약한 학습기 (실수로부터 학습):} 여전히 부족하여 교수님께 조언을 구합니다. 교수님은 "교차 검증(Cross Validation)이 답인 경우가 많아!"라고 조언합니다. 이 역시 이전 규칙들로 틀린 문제에 집중합니다.
    \item \textbf{최종 결합:} 이제 세 가지 규칙("D 옵션 피하기", "과대적합 선택", "교차 검증 선택")을 각각의 중요도에 따라 가중치(Weights)를 부여하여 결합합니다. 예를 들어, "D 옵션 피하기"는 가중치를 낮게, "교차 검증 선택"은 가중치를 높게 부여하는 식입니다.
    \item \textbf{결과:} 이 세 가지 약한 규칙을 적절히 결합한 결과, 최종 시험에서 90점을 받아 A학점을 받습니다.
\end{enumerate}
\end{tcolorbox}

이 비유에서 각 조언은 \textbf{약한 학습기(Weak Learner)}에 해당합니다. 이들은 단독으로는 완벽하지 않지만, \textbf{이전 모델의 실수로부터 학습}하여 \textbf{순차적으로(Iteratively)} \textbf{가중치를 부여하여 결합(Additively Combine)}함으로써 강력한 \textbf{강한 학습기(Strong Learner)}를 만들어냅니다.

\newpage
\section{그래디언트 부스팅(Gradient Boosting)의 작동 원리}

\subsection{핵심 아이디어: 약한 학습기의 점진적 결합}
그래디언트 부스팅의 핵심은 다음과 같습니다.
\begin{itemize}
    \item \textbf{약한 학습기(Simple Models):} 깊이가 얕은 결정 트리(예: 스텀프, Stump)와 같은 단순한 모델들을 사용합니다.
    \item \textbf{가산적 결합(Additively Combine):} 이 약한 학습기들을 하나씩 추가하면서 이전 모델의 예측을 보완합니다.
    \item \textbf{반복적 학습(Iteratively):} 다음 약한 학습기는 이전까지의 앙상블 모델이 저지른 실수(잔차)를 학습하여 보정합니다.
\end{itemize}
최종 모델은 각 약한 학습기에 \textbf{가중치(\texttt{\lambda})}를 곱하여 합산한 형태가 됩니다:
\[ F(X) = \sum_{m=1}^{M} \lambda_m T_m(X) \]
여기서 $T_m(X)$는 $m$번째 약한 학습기이고, $\lambda_m$은 해당 학습기의 가중치입니다.

\subsection{회귀 문제 예시를 통한 단계별 설명}
다음은 회귀 문제에서 그래디언트 부스팅이 어떻게 작동하는지 구체적인 예시를 통해 설명합니다.

\begin{tcolorbox}[example, title={회귀 문제 예시}] % 마스터 템플릿의 example 스타일 사용
우리는 예측 변수(X)와 반응 변수(Y)를 가진 데이터를 사용하여 회귀 모델을 구축하고자 합니다.

\begin{enumerate}
    \item \textbf{1단계: 첫 번째 약한 학습기 학습 (\texttt{T0})}
    \begin{itemize}
        \item 가장 먼저, 원본 데이터(X, Y)에 대해 첫 번째 약한 학습기($T_0$)를 학습시킵니다.
        \item 약한 학습기는 보통 깊이가 1인 스텀프(Stump)입니다. 즉, 데이터를 한 번만 분할합니다.
        \item 이 스텀프는 가중 평균 제곱 오차(Weighted MSE)를 최소화하는 분할 지점을 찾습니다.
        \item 각 분할된 영역에서는 해당 영역의 Y 값들의 평균을 예측값으로 사용합니다.
        \item 예를 들어, X가 6.8보다 작으면 -0.01, X가 6.8보다 크면 7.5를 예측합니다.
    \end{itemize}

    \item \textbf{2단계: 잔차(Residuals) 계산 (\texttt{R0})}
    \begin{itemize}
        \item 첫 번째 모델($T_0$)의 예측값($\hat{Y}_0$)과 실제 값(Y) 사이의 차이인 잔차($R_0 = Y - \hat{Y}_0$)를 계산합니다.
        \item 이 잔차는 $T_0$가 얼마나 틀렸는지를 나타냅니다.
    \end{itemize}

    \item \textbf{3단계: 잔차에 대한 두 번째 약한 학습기 학습 (\texttt{T1})}
    \begin{itemize}
        \item 이제 \textbf{잔차($R_0$)}를 새로운 목표 변수(Y)로 간주하고, 두 번째 약한 학습기($T_1$)를 학습시킵니다.
        \item $T_1$은 $R_0$를 최소화하는 분할 지점을 찾고, 각 영역에서 $R_0$의 평균을 예측합니다.
        \item 예를 들어, X가 3.5보다 작으면 잔차를 2로, X가 3.5보다 크면 잔차를 -1.5로 예측합니다.
        \item 이 단계는 이전 모델의 \textbf{실수(잔차)를 수정하는 방법}을 학습하는 것입니다.
    \end{itemize}

    \item \textbf{4단계: 학습률(\texttt{\lambda})을 이용한 모델 결합}
    \begin{itemize}
        \item 현재까지의 최종 모델($F_1$)은 첫 번째 모델($T_0$)과 두 번째 모델($T_1$)의 예측을 결합하여 만듭니다.
        \item 이때, 두 번째 모델($T_1$)의 예측에 \textbf{학습률(\texttt{\lambda})}을 곱하여 더합니다.
        \[ F_1(X) = T_0(X) + \lambda \cdot T_1(X) \]
        \item 강의에서는 \texttt{\lambda = 0.5}로 설정했습니다. 이 학습률은 모델이 잔차를 얼마나 강하게 반영할지 조절하는 역할을 합니다. 너무 크게 설정하면 과대적합될 위험이 있습니다.
    \end{itemize}

    \item \textbf{5단계: 반복 (다음 잔차 계산 및 학습)}
    \begin{itemize}
        \item 이제 새로운 최종 모델($F_1$)에 대한 잔차($R_1 = Y - F_1(X)$)를 계산합니다.
        \item 이 $R_1$을 목표 변수로 하여 세 번째 약한 학습기($T_2$)를 학습시킵니다.
        \item 이 과정을 미리 정해진 반복 횟수(예: 100번) 또는 잔차 개선이 미미해질 때까지 반복합니다.
        \item 최종 모델은 모든 약한 학습기들의 예측을 학습률과 함께 합산한 형태가 됩니다.
        \[ F_M(X) = T_0(X) + \lambda \cdot T_1(X) + \lambda \cdot T_2(X) + \dots + \lambda \cdot T_M(X) \]
    \end{itemize}
\end{enumerate}
\end{tcolorbox}

\subsection{공식화된 그래디언트 부스팅 절차}
위의 예시를 일반화하면 다음과 같은 절차로 요약할 수 있습니다.

\begin{enumerate}
    \item \textbf{초기 모델 학습:} 훈련 데이터에 대해 첫 번째 약한 모델 $T_0$를 학습합니다. (예: 모든 Y 값의 평균을 예측하는 모델)
    \item \textbf{반복:} $m=1$부터 미리 정해진 횟수 또는 종료 조건이 충족될 때까지 다음을 반복합니다.
    \begin{enumerate}
        \item \textbf{잔차 계산:} 현재 앙상블 모델 $F_{m-1}(X)$의 잔차 $R_{m-1} = Y - F_{m-1}(X)$를 계산합니다.
        \item \textbf{약한 학습기 학습:} 계산된 잔차 $R_{m-1}$를 목표 변수로 사용하여 새로운 약한 학습기 $T_m(X)$를 학습합니다.
        \item \textbf{모델 업데이트:} 현재 앙상블 모델을 업데이트합니다.
        \[ F_m(X) = F_{m-1}(X) + \lambda \cdot T_m(X) \]
        여기서 $\lambda$는 \textbf{학습률(Learning Rate)}입니다.
    \end{enumerate}
    \item \textbf{최종 예측:} 최종 앙상블 모델 $F_M(X)$를 사용하여 예측을 수행합니다.
\end{enumerate}

\subsection{학습률(\texttt{\lambda})의 역할과 중요성}
학습률 $\lambda$는 각 단계에서 새로운 약한 학습기의 기여도를 조절하는 중요한 하이퍼파라미터입니다.

\begin{itemize}
    \item \textbf{과대적합 방지:} $\lambda$가 너무 크면, 새로운 약한 학습기가 잔차를 너무 공격적으로 수정하여 모델이 훈련 데이터에 과대적합될 수 있습니다. 마치 골프 게임에서 공을 너무 세게 쳐서 목표를 지나쳐 버리는 것과 같습니다.
    \item \textbf{안정적인 학습:} $\lambda$를 작게 설정하면, 모델은 잔차를 조금씩, 점진적으로 수정하며 더 안정적으로 학습합니다. 이는 수렴 속도는 느리지만, 더 견고한 모델을 만들 가능성이 높습니다.
    \item \textbf{잔차 감소 보장:} $\lambda$를 적절히 설정하면, 각 반복마다 잔차가 감소하는 것을 보장할 수 있습니다.
\end{itemize}

\begin{tcolorbox}[importantbox, title={학습률(\texttt{\lambda})의 중요성}] % 마스터 템플릿의 importantbox 스타일 사용
학습률(\texttt{\lambda})은 그래디언트 부스팅의 성능에 매우 큰 영향을 미칩니다. 너무 크면 과대적합과 불안정한 학습을 초래하고, 너무 작으면 학습 시간이 매우 길어질 수 있습니다. 적절한 \texttt{\lambda} 값은 교차 검증(Cross-validation)을 통해 찾아야 합니다.
\end{tcolorbox}

\subsection{실제 구현 시 유의사항}
실제로 그래디언트 부스팅을 구현할 때는 모든 약한 학습기들을 하나의 거대한 트리로 명시적으로 결합하지 않습니다. 대신, 각 약한 학습기($T_0, T_1, \dots, T_M$)를 개별적으로 저장하고, 새로운 데이터에 대한 예측(추론) 시에는 각 트리의 예측값에 학습률을 곱하여 모두 더하는 방식으로 최종 예측값을 얻습니다.

\[ \hat{Y}_{final}(X) = T_0(X) + \lambda \cdot T_1(X) + \lambda \cdot T_2(X) + \dots + \lambda \cdot T_M(X) \]

\begin{tipbox}[머신러닝에서의 '추론(Inference)' 용어] % 마스터 템플릿의 tipbox 사용
통계학에서 '추론(Inference)'은 모집단에 대한 결론을 도출하는 것을 의미합니다. 하지만 머신러닝 분야에서는 종종 '추론'이라는 용어가 \textbf{모델이 새로운 데이터에 대해 예측(Prediction)을 수행하는 과정}을 의미하기도 합니다. 혼동을 피하기 위해 본 강의에서는 '예측'이라는 용어를 주로 사용합니다.
\end{tipbox}

\newpage
\section{그래디언트 부스팅의 수학적 정당화: 그래디언트 하강(Gradient Descent)과의 연결}

\subsection{잔차 감소의 직관적 설명}
우리는 각 단계에서 약한 학습기가 이전 모델의 잔차를 학습하도록 합니다. 이는 모델이 이전에 저지른 실수를 점진적으로 수정해 나가는 과정입니다.
\[ R_m = R_{m-1} - \lambda \cdot T_m(X) \]
여기서 $R_m$은 $m$번째 반복 후의 잔차이고, $R_{m-1}$은 이전 잔차입니다. $T_m(X)$는 이전 잔차를 예측하는 모델입니다. 만약 $T_m(X)$가 잔차를 잘 예측한다면, 새로운 잔차 $R_m$은 이전 잔차 $R_{m-1}$보다 작아질 것입니다. 즉, 모델의 오차가 점진적으로 감소합니다.

\subsection{최적화의 기본 원리: 그래디언트 하강(Gradient Descent)}
어떤 함수(예: 손실 함수, Loss Function)를 최소화하고자 할 때, 우리는 보통 해당 함수의 도함수(Derivative) 또는 편도함수(Partial Derivative)를 계산하여 0으로 설정함으로써 최솟값을 찾습니다. 하지만 복잡한 함수에서는 이를 분석적으로 풀기 어렵습니다. 이때 \textbf{그래디언트 하강(Gradient Descent)} 알고리즘이 유용하게 사용됩니다.

\begin{tcolorbox}[summary, title={그래디언트 하강(Gradient Descent)의 핵심}] % 마스터 템플릿의 summary 스타일 사용
그래디언트 하강은 함수의 최솟값을 찾기 위한 \textbf{반복적인(Iterative)} 최적화 알고리즘입니다.
\begin{itemize}
    \item \textbf{목표:} 손실 함수 $L(W)$를 최소화하는 파라미터 $W$를 찾습니다.
    \item \textbf{절차:}
    \begin{enumerate}
        \item \textbf{초기화:} 파라미터 $W$를 임의의 값으로 초기화합니다.
        \item \textbf{기울기 계산:} 현재 파라미터 $W$에서 손실 함수 $L(W)$의 기울기(Gradient, $\nabla L(W)$)를 계산합니다. 기울기는 함수가 가장 빠르게 증가하는 방향을 나타냅니다.
        \item \textbf{파라미터 업데이트:} 기울기의 \textbf{반대 방향}으로 파라미터를 업데이트합니다. 이때 \textbf{학습률(\texttt{\eta})}을 곱하여 이동할 스텝의 크기를 조절합니다.
        \[ W_{new} = W_{old} - \eta \cdot \nabla L(W_{old}) \]
        \item \textbf{반복:} 수렴하거나 미리 정해진 반복 횟수에 도달할 때까지 2, 3단계를 반복합니다.
    \end{enumerate}
\end{itemize}
\end{tcolorbox}

\begin{tcolorbox}[example, title={그래디언트 하강 비유: 어두운 산에서 하산하기}] % 마스터 템플릿의 example 스타일 사용
어두운 밤, 산 정상에 홀로 서서 가장 낮은 곳(음식과 물이 있는 곳)으로 내려가야 한다고 상상해 보세요. 아무것도 보이지 않지만, 발밑의 경사(기울기)는 느낄 수 있습니다.
\begin{itemize}
    \item \textbf{기울기 측정:} 발밑에서 가장 가파른 내리막길 방향을 찾습니다. 이것이 기울기의 반대 방향입니다.
    \item \textbf{스텝 크기 조절:} 경사가 가파르면(기울기 절댓값이 크면) 큰 걸음으로 빠르게 내려가고, 경사가 완만해지면(기울기 절댓값이 작아지면) 작은 걸음으로 조심스럽게 내려갑니다. 이것이 학습률(\texttt{\eta})과 기울기 크기의 곱으로 스텝 크기를 조절하는 원리입니다.
    \item \textbf{반복:} 계속해서 발밑의 경사를 느끼고, 가장 가파른 내리막길 방향으로 걸음을 옮기는 과정을 반복합니다. 결국 가장 낮은 곳(최솟값)에 도달하게 됩니다.
\end{itemize}
\end{tcolorbox}

\begin{figure}[h!]
    \centering
    % \includegraphics[width=0.7\textwidth]{example-gradient-descent.png}  % Image not found: example-gradient-descent.png % Placeholder for actual image
    \caption{그래디언트 하강의 작동 원리 (가상의 이미지)}
    \label{fig:gradient_descent}
\end{figure}

\subsubsection{학습률(\texttt{\eta})의 영향}
\begin{itemize}
    \item \textbf{\texttt{\eta}가 너무 작을 때:} 매우 작은 걸음으로 이동하므로 최솟값에 도달하는 데 시간이 오래 걸립니다.
    \item \textbf{\texttt{\eta}가 너무 클 때:} 너무 큰 걸음으로 이동하여 최솟값을 건너뛰거나, 오히려 발산하여 최솟값에서 멀어질 수 있습니다. 마치 산을 내려가다가 너무 크게 점프해서 반대편 언덕으로 넘어가 버리는 것과 같습니다.
\end{itemize}

\subsubsection{볼록 함수와 지역 최솟값}
\begin{itemize}
    \item \textbf{볼록 함수(Convex Function):} 손실 함수가 볼록하면 그래디언트 하강은 항상 전역 최솟값(Global Minimum)을 찾을 수 있습니다. 로지스틱 회귀(Logistic Regression)의 손실 함수는 볼록하지만, 분석적으로 해를 찾기 어렵기 때문에 그래디언트 하강을 사용합니다.
    \item \textbf{비볼록 함수(Non-convex Function):} 손실 함수가 비볼록하면 그래디언트 하강은 지역 최솟값(Local Minimum)에 갇힐 수 있습니다.
        \begin{itemize}
            \item \textbf{해결책:}
                \begin{itemize}
                    \item \textbf{무작위 재시작(Random Restarts):} 여러 시작점에서 그래디언트 하강을 수행하여 가장 좋은 최솟값을 선택합니다.
                    \item \textbf{확률적 그래디언트 하강(Stochastic Gradient Descent, SGD):} 데이터의 일부만 사용하여 기울기를 계산하여 지역 최솟값에서 벗어날 가능성을 높입니다 (CS109B에서 다룸).
                    \item \textbf{실용적 접근:} 지역 최솟값이 전역 최솟값만큼 좋지 않더라도, 교차 검증을 통해 충분히 좋은 성능을 보인다면 만족하는 경우도 많습니다. 딥러닝 분야에서는 지역 최솟값이 전역 최솟값과 크게 다르지 않다는 연구 결과도 있습니다.
                \end{itemize}
        \end{itemize}
\end{itemize}

\subsection{부스팅을 그래디언트 하강으로 해석}
이제 그래디언트 부스팅을 그래디언트 하강의 관점에서 이해해 봅시다.

\subsubsection{예측 공간에서의 최적화}
일반적인 그래디언트 하강은 모델의 파라미터(예: 선형 회귀의 $\beta$ 값, 신경망의 가중치 $W$)를 최적화합니다. 하지만 결정 트리와 같은 \textbf{비모수 모델(Non-parametric Models)}은 고정된 수의 파라미터로 정의되지 않습니다. (트리는 데이터를 분할할 때마다 새로운 노드가 생기므로 파라미터 수가 고정되어 있지 않습니다.)

따라서 그래디언트 부스팅은 파라미터 공간이 아닌 \textbf{함수 공간(Function Space)} 또는 \textbf{예측 공간(Prediction Space)}에서 최적화를 수행하는 것으로 볼 수 있습니다.

\begin{itemize}
    \item \textbf{손실 함수:} 평균 제곱 오차(MSE)를 손실 함수로 사용한다고 가정해 봅시다.
    \[ L(Y, \hat{Y}) = \frac{1}{N} \sum_{i=1}^{N} (Y_i - \hat{Y}_i)^2 \]
    \item \textbf{기울기 계산:} 이 손실 함수를 예측값 $\hat{Y}_i$에 대해 미분하면 다음과 같습니다.
    \[ \frac{\partial L}{\partial \hat{Y}_i} = -2(Y_i - \hat{Y}_i) \]
    \item \textbf{잔차와의 관계:} 이 기울기는 실제 값과 예측값의 차이, 즉 \textbf{잔차(Residual)}에 비례합니다 (부호는 반대).
    \[ -\frac{1}{2} \frac{\partial L}{\partial \hat{Y}_i} = Y_i - \hat{Y}_i = R_i \]
    따라서, 손실 함수를 최소화하는 가장 이상적인 방향은 \textbf{잔차의 방향}입니다. 만약 $\hat{Y}_i$를 파라미터로 직접 최적화한다면, 한 번의 스텝으로 $\hat{Y}_i = Y_i$가 되어 완벽하게 과대적합됩니다. 이는 모델이 아무것도 학습하지 않은 것과 같습니다.
\end{itemize}

\subsubsection{약한 학습기가 기울기를 근사하는 역할}
그래디언트 부스팅은 이 문제를 다음과 같이 해결합니다.
\begin{itemize}
    \item \textbf{이상적인 방향:} 각 데이터 포인트에 대해 손실 함수의 음의 기울기(Negative Gradient)는 해당 데이터 포인트의 잔차와 같습니다. 이것이 모델이 나아가야 할 '이상적인' 방향입니다.
    \item \textbf{약한 학습기의 근사:} 하지만 우리는 $\hat{Y}_i$를 직접 업데이트하는 대신, \textbf{약한 학습기($T_m$)}를 사용하여 이 잔차(즉, 음의 기울기)를 \textbf{근사(Approximate)}합니다.
    \item \textbf{점진적 개선:} $T_m$은 매우 단순한 모델(스텀프)이므로, 잔차를 완벽하게 근사하지 못합니다. 하지만 이 불완전한 근사를 통해 모델은 잔차를 조금씩, 점진적으로 수정해 나갑니다.
    \item \textbf{학습률(\texttt{\lambda})의 재해석:} 이때 학습률 $\lambda$는 그래디언트 하강에서 스텝 크기를 조절하는 $\eta$와 동일한 역할을 합니다. 즉, 약한 학습기가 근사한 기울기 방향으로 얼마나 크게 스텝을 밟을지 결정합니다.
\end{itemize}

\begin{figure}[h!]
    \centering
    % \includegraphics[width=0.7\textwidth]{example-gradient-boosting-path.png}  % Image not found: example-gradient-boosting-path.png % Placeholder for actual image
    \caption{그래디언트 부스팅의 함수 공간에서의 경로 (가상의 이미지)}
    \label{fig:gradient_boosting_path}
\end{figure}

위 그림 \ref{fig:gradient_boosting_path}에서 볼 수 있듯이, 이상적인 경로는 잔차 방향으로 한 번에 최솟값에 도달하는 것이지만 (과대적합), 그래디언트 부스팅은 약한 학습기를 통해 잔차(기울기)를 근사하고, 학습률(\texttt{\lambda})을 조절하며 점진적으로 최솟값에 가까워집니다.

\begin{tcolorbox}[summary, title={그래디언트 부스팅 = 함수 공간에서의 그래디언트 하강}] % 마스터 템플릿의 summary 스타일 사용
그래디언트 부스팅은 모델의 파라미터가 아닌 \textbf{모델 자체(함수)}를 최적화하는 그래디언트 하강의 한 형태로 볼 수 있습니다. 각 약한 학습기는 현재 앙상블 모델의 손실 함수에 대한 \textbf{음의 기울기(잔차)}를 근사하고, 학습률(\texttt{\lambda})에 따라 그 방향으로 모델을 업데이트합니다. 이 과정을 반복하여 최종적으로 손실을 최소화하는 강력한 모델을 구축합니다.
\end{tcolorbox}

\subsection{수렴 보장 및 학습률(\texttt{\lambda}) 튜닝}
\begin{itemize}
    \item \textbf{수렴 보장:} 그래디언트 부스팅은 그래디언트 하강의 원리를 따르므로, 학습률 $\lambda$가 충분히 작다면 모델이 최솟값으로 수렴하는 것을 보장할 수 있습니다.
    \item \textbf{학습률 튜닝:} $\lambda$는 교차 검증(Cross-validation)을 통해 최적의 값을 찾아야 합니다.
        \begin{itemize}
            \item \textbf{시각적 확인:} 학습 과정에서 손실(Loss)이 꾸준히 감소하거나 정확도(Accuracy)가 꾸준히 증가하는지 확인합니다.
            \item \textbf{불안정한 학습:} 만약 손실이 급격히 변동하거나 정확도가 오르내린다면, $\lambda$가 너무 높게 설정되었을 가능성이 큽니다.
            \item \textbf{고급 기법:} 109B와 같은 고급 과정에서는 기울기의 크기에 반비례하여 $\lambda$를 조절하는 Adam, Momentum과 같은 고급 최적화 기법을 배우게 됩니다.
        \end{itemize}
\end{itemize}

\subsection{종료 조건}
그래디언트 부스팅은 다음 두 가지 방법 중 하나로 학습을 종료합니다.
\begin{itemize}
    \item \textbf{최대 반복 횟수 제한:} 미리 정해진 수의 약한 학습기(예: 100개)를 생성한 후 학습을 종료합니다.
    \item \textbf{성능 개선 임계값:} 새로운 약한 학습기를 추가해도 손실 함수(예: 잔차 제곱합)의 개선이 특정 임계값(예: 0.1\%) 미만일 경우 학습을 종료합니다. 이는 불필요한 계산을 줄이고 과대적합을 방지하는 데 도움이 됩니다.
\end{itemize}

\newpage
\section{체크리스트: 학습 점검}
다음 체크리스트를 통해 그래디언트 부스팅에 대한 이해도를 점검해 보세요.

\begin{itemize}
    \item \textbf{부스팅의 동기 이해:} 왜 랜덤 포레스트 외에 부스팅이 필요한지 설명할 수 있는가?
    \item \textbf{약한 학습기 개념:} 약한 학습기가 무엇이며, 왜 부스팅에서 약한 학습기를 사용하는지 설명할 수 있는가? (예: 스텀프)
    \item \textbf{그래디언트 부스팅 절차:} 회귀 예시를 통해 그래디언트 부스팅의 단계별 작동 원리를 설명할 수 있는가? (잔차 학습, 가산적 결합)
    \item \textbf{학습률(\texttt{\lambda})의 역할:} 학습률이 무엇이며, 모델 학습에 어떤 영향을 미치는지 설명할 수 있는가?
    \item \textbf{그래디언트 하강 기본 원리:} 그래디언트 하강이 무엇이며, 어떻게 최솟값을 찾아가는지 설명할 수 있는가? (기울기, 스텝 크기)
    \item \textbf{부스팅과 그래디언트 하강의 연결:} 그래디언트 부스팅이 왜 함수 공간에서의 그래디언트 하강으로 해석될 수 있는지 설명할 수 있는가? (잔차가 음의 기울기)
    \item \textbf{비모수 모델의 이해:} 결정 트리가 왜 비모수 모델인지 설명할 수 있는가?
    \item \textbf{학습률 튜닝 및 종료 조건:} 학습률을 어떻게 튜닝하고, 언제 학습을 종료해야 하는지 설명할 수 있는가?
    \item \textbf{표 형식 데이터에서의 중요성:} 부스팅이 표 형식 데이터에서 왜 강력한지 설명할 수 있는가?
\end{itemize}

\newpage
\section{FAQ: 자주 묻는 질문}

\begin{faqbox}[title=\textbf{Q: 왜 부스팅이 랜덤 포레스트보다 튜닝하기 어려운가요?}] % faqbox 스타일 사용
\textbf{A:} 부스팅은 약한 학습기들을 순차적으로 학습시키기 때문에, 각 단계의 결정이 다음 단계에 영향을 미칩니다. 이로 인해 학습률(\texttt{\lambda}), 약한 학습기의 깊이, 반복 횟수 등 여러 하이퍼파라미터 간의 상호작용이 복잡해집니다. 반면 랜덤 포레스트는 트리를 독립적으로 학습시키므로 튜닝이 상대적으로 간단합니다. 부스팅은 과대적합에 더 취약할 수 있어 신중한 튜닝이 필요합니다.
\end{faqbox}

\begin{faqbox}[title=\textbf{Q: 약한 학습기(Weak Learner)란 정확히 무엇인가요?}] % faqbox 스타일 사용
\textbf{A:} 약한 학습기는 단독으로는 예측 성능이 좋지 않지만, 무작위 추측보다는 약간 더 나은 성능을 보이는 모델을 의미합니다. 그래디언트 부스팅에서는 주로 깊이가 1인 결정 트리인 \textbf{스텀프(Stump)}를 약한 학습기로 사용합니다. 스텀프는 데이터를 단 한 번만 분할하여 매우 단순한 규칙을 학습합니다. 이러한 단순함 덕분에 과대적합 위험이 낮고, 여러 개를 결합했을 때 강력한 시너지를 낼 수 있습니다.
\end{faqbox}

\begin{faqbox}[title=\textbf{Q: 학습률(\texttt{\lambda})이 너무 크거나 작으면 모델 학습에 어떤 문제가 발생하나요?}] % faqbox 스타일 사용
\textbf{A:}
\begin{itemize}
    \item \textbf{\texttt{\lambda}가 너무 크면:} 모델이 잔차를 너무 공격적으로 수정하여 최솟값을 건너뛰거나 발산할 수 있습니다. 이는 학습이 불안정해지고 과대적합될 가능성을 높입니다.
    \item \textbf{\texttt{\lambda}가 너무 작으면:} 모델이 매우 작은 스텝으로 이동하므로 최솟값에 도달하는 데 매우 오랜 시간이 걸립니다. 이는 계산 비용이 커지고 학습 효율이 떨어집니다.
\end{itemize}
따라서 적절한 \texttt{\lambda} 값을 찾는 것이 중요하며, 일반적으로 교차 검증을 통해 최적의 값을 탐색합니다.
\end{faqbox}

\begin{faqbox}[title=\textbf{Q: 결정 트리가 왜 비모수 모델(Non-parametric Model)인가요?}] % faqbox 스타일 사용
\textbf{A:} 비모수 모델은 데이터의 특정 분포를 가정하지 않거나, 모델의 복잡도(즉, 파라미터의 수)가 데이터에 따라 유연하게 변하는 모델을 의미합니다. 결정 트리는 데이터를 분할할 때마다 새로운 노드와 분할 규칙이 생성되므로, 모델의 구조와 파라미터 수가 고정되어 있지 않습니다. 반면 선형 회귀나 로지스틱 회귀는 $\beta_0, \beta_1, \dots$와 같이 고정된 수의 파라미터를 가집니다. 이러한 유연성 때문에 결정 트리는 비모수 모델로 분류됩니다.
\end{faqbox}

\begin{faqbox}[title=\textbf{Q: 머신러닝에서 '추론(Inference)'이라는 용어는 통계적 추론과 어떻게 다른가요?}] % faqbox 스타일 사용
\textbf{A:}
\begin{itemize}
    \item \textbf{통계적 추론:} 모집단 특성(예: 평균, 분산)에 대해 표본 데이터를 기반으로 결론을 도출하거나 가설을 검정하는 과정입니다. 불확실성을 정량화하고 신뢰 구간을 제시하는 것이 중요합니다.
    \item \textbf{머신러닝 추론:} 훈련된 머신러닝 모델이 새로운, 보지 못한 데이터에 대해 예측(Prediction)을 수행하는 과정을 의미합니다. 주로 모델의 예측 성능에 초점을 맞춥니다.
\end{itemize}
강의에서는 혼동을 피하기 위해 머신러닝 맥락에서는 '예측'이라는 용어를 사용하는 것이 좋다고 강조합니다.
\end{faqbox}

\newpage
\section{빠르게 훑어보기 (1페이지 요약)}

\begin{tcolorbox}[summary, title=\textbf{부스팅 핵심 요약}] % summary 스타일 사용
\textbf{목표:} 약한 학습기(Weak Learners)를 순차적으로 결합하여 모델의 \textbf{편향(Bias)}을 줄이고 강력한 모델을 구축.

\textbf{랜덤 포레스트와의 차이:}
\begin{itemize}
    \item \textbf{랜덤 포레스트:} 독립적인 트리 병렬 학습, \textbf{분산(Variance)} 감소.
    \item \textbf{부스팅:} 순차적인 트리 학습, 이전 모델의 \textbf{실수(잔차)} 보완, \textbf{편향} 감소.
\end{itemize}
\end{tcolorbox}

\begin{tcolorbox}[infobox, title=\textbf{그래디언트 부스팅 작동 원리}] % infobox 스타일 사용
\textbf{1. 초기 모델:} 원본 데이터에 첫 번째 약한 학습기($T_0$) 학습.
\textbf{2. 잔차 계산:} $T_0$의 예측 오차(잔차 $R_0$) 계산.
\textbf{3. 잔차 학습:} $R_0$를 목표로 두 번째 약한 학습기($T_1$) 학습.
\textbf{4. 모델 결합:} $F_1(X) = T_0(X) + \lambda \cdot T_1(X)$ (학습률 $\lambda$ 사용).
\textbf{5. 반복:} 새로운 $F_1(X)$의 잔차를 계산하고, 이를 학습하는 다음 약한 학습기를 추가하는 과정 반복.
\textbf{핵심:} 다음 모델은 이전 모델의 \textbf{실수(잔차)}로부터 학습하여 점진적으로 개선.
\end{tcolorbox}

\begin{tcolorbox}[warningbox, title=\textbf{그래디언트 하강(Gradient Descent)과의 연결}] % warningbox 스타일 사용
\textbf{1. 최적화 목표:} 손실 함수 $L(Y, \hat{Y})$를 최소화.
\textbf{2. 기울기(Gradient):} 손실 함수를 예측값 $\hat{Y}$에 대해 미분하면 \textbf{잔차($Y - \hat{Y}$)}에 비례.
\textbf{3. 함수 공간 최적화:} 결정 트리는 비모수 모델이므로 파라미터 공간이 아닌 \textbf{함수 공간}에서 최적화.
\textbf{4. 약한 학습기의 역할:} 각 약한 학습기($T_m$)는 현재 앙상블 모델의 손실 함수에 대한 \textbf{음의 기울기(잔차)}를 근사.
\textbf{5. 학습률(\texttt{\lambda}):} 그래디언트 하강의 스텝 크기(\texttt{\eta})와 동일하게, 기울기 방향으로 얼마나 크게 이동할지 조절.
\textbf{결론:} 그래디언트 부스팅은 함수 공간에서 \textbf{잔차(음의 기울기)}를 약한 학습기로 근사하며 점진적으로 최솟값을 찾아가는 그래디언트 하강의 한 형태.
\end{tcolorbox}

\begin{tcolorbox}[importantbox, title=\textbf{중요 사항}] % importantbox 스타일 사용
\textbf{표 형식 데이터:} 부스팅과 랜덤 포레스트는 표 형식 데이터에서 신경망보다 우수한 성능을 보이는 경우가 많음.
\textbf{학습률(\texttt{\lambda}) 튜닝:} 과대적합 방지 및 안정적인 학습을 위해 교차 검증을 통한 신중한 튜닝 필수.
\textbf{종료 조건:} 최대 반복 횟수 또는 손실 개선 임계값 설정.
\end{tcolorbox}

\end{document}
